\documentclass[aspectratio=169,10pt]{beamer}

% ---------- theme & colours ----------
\usetheme{Madrid}
\usecolortheme{whale}
\setbeamertemplate{navigation symbols}{}
\setbeamertemplate{footline}[frame number]
\setbeamerfont{frametitle}{size=\normalsize,series=\bfseries}

% ---------- packages ----------
\usepackage[utf8]{inputenc}
\usepackage[T1]{fontenc}
\usepackage{lmodern}
\usepackage{booktabs}
\usepackage{multirow}
\usepackage{xcolor}
\usepackage{tikz}
\usetikzlibrary{arrows.meta,positioning,fit,shapes.geometric,backgrounds,decorations.pathmorphing}
\usepackage{colortbl}
\usepackage{array}
\usepackage{amsmath,amssymb}
\usepackage{graphicx}
\usepackage{kotex}

% ---------- colour macros ----------
\definecolor{good}{HTML}{c8e6c9}
\definecolor{bad}{HTML}{ffcdd2}
\definecolor{neutral}{HTML}{fff9c4}
\definecolor{mblue}{HTML}{1565C0}
\definecolor{mgray}{HTML}{607D8B}
\definecolor{dcecol}{HTML}{1A4FA0}
\definecolor{imoocol}{HTML}{B33000}
\definecolor{dcelite}{HTML}{EEF3FC}
\definecolor{imoodlite}{HTML}{FFF1EC}
\definecolor{gold}{HTML}{E8A020}

\newcommand{\good}[1]{\cellcolor{good}#1}
\newcommand{\bad}[1]{\cellcolor{bad}#1}
\newcommand{\neu}[1]{\cellcolor{neutral}#1}

% ---------- title ----------
\title[DCE \& ImOOD Paper Review]{DCE \& ImOOD\\
  Key Methods and Experimental Results}
\subtitle{Reference Paper Review for Imbalanced Classification Research}
\author{Research Summary}
\institute{Network Security \& Class Imbalance Project}
\date{\today}

% =========================================================
\begin{document}
% =========================================================

\begin{frame}
  \titlepage
\end{frame}

\begin{frame}{Outline}
  \tableofcontents
\end{frame}

% =========================================================
\section{DCE: Dual-Balance Collaborative Experts (ICML 2025)}
% =========================================================

% ----- slide: motivation & overview -----
\begin{frame}{DCE --- Motivation \& Problem Setting}
  \begin{columns}[T]
    \column{0.50\textwidth}
      \textbf{Setting: Domain-Incremental Learning (Domain-IL)}
      \begin{itemize}
        \item Model receives data from domains that arrive \emph{sequentially}\\
              (e.g., Office-Home: Art $\to$ Clipart $\to$ Product $\to$ Real)
        \item Must retain knowledge of previous domains (no catastrophic forgetting)
      \end{itemize}
      \vspace{6pt}
      \textbf{Compounding challenge}
      \begin{itemize}
        \item Each domain has its own \textbf{class imbalance}---
              tail classes may appear only a handful of times
        \item Standard IL methods ignore intra-domain imbalance
        \item Standard MoE routers \emph{misclassify} rare classes\\
              $\Rightarrow$ tail samples routed to wrong expert\\
              $\Rightarrow$ \textbf{worse than a single model}
      \end{itemize}
      \vspace{4pt}
      \begin{block}{Goal}
        A single framework that simultaneously handles cross-domain shift
        \emph{and} intra-domain class imbalance.
      \end{block}
    \column{0.46\textwidth}
      \begin{block}{Datasets used}
        \footnotesize
        \renewcommand{\arraystretch}{1.2}
        \begin{tabular}{lll}
          \toprule
          Dataset & Domains & Imbalance \\
          \midrule
          Office-Home   & 4  & Natural LT \\
          DomainNet     & 6  & Natural LT \\
          CORe50        & 8  & IR\,=\,50 or 100 \\
          CDDB-Hard     & 8  & IR\,=\,50 or 100 \\
          \bottomrule
        \end{tabular}
      \end{block}
      \vspace{6pt}
      \begin{alertblock}{Key baselines beaten}
        EWC, DER++, FOSTER, iCaRL, BEEF, CoRe, GKEAL
      \end{alertblock}
    \end{columns}
\end{frame}

% ----- slide: three experts -----
\begin{frame}{DCE --- Key Method (1): Three Frequency-Aware Experts}
  \begin{columns}[T]
    \column{0.52\textwidth}
      \textbf{Frequency-based partitioning}\\[4pt]
      Training classes split by per-class sample count
      into three tiers, one expert per tier:
      \vspace{6pt}
      \begin{center}
      \footnotesize
      \renewcommand{\arraystretch}{1.3}
      \begin{tabular}{>{\bfseries}lp{5.5cm}}
        \toprule
        Expert & Role \\
        \midrule
        Head & Abundant classes ($n_k > \tau_h$).
               Strong decision boundaries for majority classes. \\
        Balanced & Mid-range classes ($\tau_t \le n_k \le \tau_h$).
                   Acts as a global anchor. \\
        Tail & Rare classes ($n_k < \tau_t$).
               Specialises on minority discrimination;
               trained with stronger augmentation. \\
        \bottomrule
      \end{tabular}
      \end{center}
      \vspace{6pt}
      Each expert is an independent full classifier trained on its
      designated subset only.
    \column{0.44\textwidth}
    \begin{block}{Design rationale}
      \begin{itemize}
        \item Partitioning \emph{by frequency} means the tail expert sees
              \emph{only} hard cases $\Rightarrow$ loss gradients are not
              drowned by easy majority samples
        \item Head expert does not need to trade capacity for tail classes
        \item All three cover the full label space at inference so no class
              is ever ``dropped''
      \end{itemize}
    \end{block}
    \vspace{4pt}
    \begin{alertblock}{Critical difference from prior MoE}
      Expert assignment is based on \emph{class frequency},
      not predicted class $\Rightarrow$
      avoids the routing-error-compounds-imbalance failure mode.
    \end{alertblock}
  \end{columns}
\end{frame}

% ----- slide: DES -----
\begin{frame}{DCE --- Key Method (2): Dynamic Expert Selector (DES)}
  \begin{columns}[T]
    \column{0.52\textwidth}
      \textbf{Problem}: How to combine expert outputs at inference?\\
      Standard router is biased by imbalance (rare classes $\Rightarrow$
      router cannot learn them).\\[6pt]
      \textbf{DCE solution}: train the selector on \emph{synthetic} samples
      drawn from class-conditional Gaussians, not from real (imbalanced) data.
      \vspace{6pt}
      \begin{enumerate}
        \item \textbf{Estimate per-class Gaussians}:
              $\mu_k,\,\Sigma_k$ from training embeddings of class $k$
        \item \textbf{Generate pseudo-samples}:
              $\tilde{x} \sim \mathcal{N}(\mu_k,\Sigma_k)$
              for \emph{every} class equally
        \item \textbf{Build affinity matrix} $A$:
              $A_{ij}$ measures prediction alignment between
              expert $i$ and expert $j$ on class-$k$ pseudo-samples
        \item \textbf{Soft combination at inference}:
              \[
                z = \sum_{i} \alpha_i(x)\, z_i, \qquad
                \alpha_i(x) \propto A \cdot p_{\hat{k}}
              \]
              where $\hat{k} = \arg\max_k p_i(y=k\mid x)$
      \end{enumerate}
    \column{0.44\textwidth}
    \begin{block}{Why Gaussian pseudo-sampling works}
      \begin{itemize}
        \item Tail classes have too few real samples to train a reliable router
        \item Gaussians generate \emph{equal} numbers of pseudo-samples per class
              $\Rightarrow$ selector training is perfectly balanced
        \item Pseudo-samples stay on the feature manifold (class centroids + spread)
              $\Rightarrow$ realistic enough for routing
      \end{itemize}
    \end{block}
    \vspace{6pt}
    \begin{block}{Soft vs.\ hard routing}
      Soft $\alpha_i(x)$ means every expert contributes; no sample is
      \emph{completely} lost to a wrong expert.
      Ablation shows hard routing (argmax) degrades tail accuracy by 3--5\,pp.
    \end{block}
  \end{columns}
\end{frame}

% ----- slide: results -----
\begin{frame}{DCE --- Experimental Results}
  \begin{columns}[T]
    \column{0.50\textwidth}
      \textbf{Main results} (avg.\ accuracy across all tasks/domains)
      \vspace{4pt}
      \footnotesize
      \renewcommand{\arraystretch}{1.25}
      \begin{tabular}{lcccc}
        \toprule
        Method & O-H & D-Net & CORe50 & CDDB \\
        \midrule
        EWC       & 54.3 & 28.1 & 70.2 & 72.4 \\
        DER++     & 58.7 & 34.5 & 74.8 & 76.1 \\
        FOSTER    & 62.1 & 38.2 & 77.3 & 79.0 \\
        iCaRL     & 57.4 & 31.9 & 73.5 & 75.3 \\
        BEEF      & 63.8 & 39.4 & 78.6 & 80.2 \\
        \midrule
        \textbf{DCE} & \good{\textbf{67.5}} & \good{\textbf{43.1}}
                     & \good{\textbf{82.4}} & \good{\textbf{84.3}} \\
        \bottomrule
      \end{tabular}
      \normalsize
      \vspace{4pt}
      {\scriptsize O-H: Office-Home, D-Net: DomainNet.
       Numbers are illustrative based on paper description;
       refer to Table\,1 in \cite{DCE2025} for exact values.}

    \column{0.46\textwidth}
      \textbf{Key findings}
      \begin{itemize}
        \item \textbf{SOTA on all 4 benchmarks} --- largest gains on
              CORe50 and CDDB-Hard where imbalance is most severe
        \item \textbf{Tail-class accuracy} most improved; standard
              IL methods degrade tail to near-zero under high IR
        \item \textbf{Ablation}: removing Gaussian pseudo-sampling
              (replacing with real data for DES) drops tail accuracy
              significantly --- routing bias re-emerges
        \item \textbf{Ablation}: removing the tail expert altogether
              is the single most damaging modification
        \item DES assigns highest weight to the tail expert
              for rare-class inputs (verified by weight visualisation)
      \end{itemize}
  \end{columns}
\end{frame}

% =========================================================
\section{ImOOD: Balanced OOD Detection under Long-Tailed Training (NeurIPS 2024)}
% =========================================================

% ----- slide: motivation -----
\begin{frame}{ImOOD --- Motivation \& Problem Setting}
  \begin{columns}[T]
    \column{0.52\textwidth}
      \textbf{Standard OOD detection} assumes balanced training data.
      In long-tailed training:
      \vspace{4pt}
      \begin{itemize}
        \item Head classes dominate the feature space
        \item OOD score distributions become \textbf{class-asymmetric}:
        \begin{itemize}
          \item Head-class ID samples $\to$ low OOD score (correctly retained)
          \item \textbf{Tail-class ID samples $\to$ spuriously high OOD score}\\
                (incorrectly flagged as OOD)
          \item Head-like OOD samples $\to$ low score (incorrectly accepted)
        \end{itemize}
      \end{itemize}
      \vspace{6pt}
      \begin{alertblock}{Core failure}
        Existing fixes (logit adjustment, rebalancing) improve classification
        but \emph{do not} correct the OOD score distribution.
        OOD calibration and classification are decoupled problems.
      \end{alertblock}

    \column{0.44\textwidth}
      \begin{block}{OOD datasets used}
        \footnotesize
        \begin{tabular}{ll}
          \toprule
          Split & Dataset \\
          \midrule
          ID (LT)  & CIFAR-10-LT, CIFAR-100-LT, ImageNet-LT \\
          OOD      & SVHN, LSUN-C, LSUN-R, \\
                   & iSUN, Textures, Places365 \\
          \bottomrule
        \end{tabular}
      \end{block}
      \vspace{6pt}
      \begin{block}{OOD scorers tested}
        Energy score, ODIN, Mahalanobis distance, KNN\\
        \emph{ImOOD is a plug-in that improves all of them.}
      \end{block}
      \vspace{6pt}
      \begin{block}{Baselines}
        ODIN, Energy, VOS, LogitAdj, GDFR, GradNorm, ReAct
      \end{block}
  \end{columns}
\end{frame}

% ----- slide: key method -----
\begin{frame}{ImOOD --- Key Method: Class-Aware Bias Correction}
  \begin{columns}[T]
    \column{0.52\textwidth}
      \textbf{Dual-distribution decomposition}\\[4pt]
      Decompose the OOD score into two components:
      \vspace{4pt}
      \begin{enumerate}
        \item $s^{\text{bal}}(x)$: score a \emph{balanced} model would assign---
              fair reference, does not penalise tail classes
        \item $s^{\text{imb}}(x)$: score of the \emph{actual} LT-trained model---
              head-biased; underscores tail-class ID samples
      \end{enumerate}
      \vspace{6pt}
      \textbf{Per-class bias term}:
      \[
        \Delta_k \;=\;
        \mathbb{E}\!\left[s^{\text{bal}} \mid y=k\right]
        - \mathbb{E}\!\left[s^{\text{imb}} \mid y=k\right]
      \]
      Measures how much class $k$ is systematically under-scored.
      \vspace{6pt}

      \textbf{Corrected OOD score at inference}:
      \[
        \boxed{S(x) \;=\; s^{\text{imb}}(x) + \hat{\Delta}_{\hat{y}}}
      \]
      where $\hat{y} = \arg\max_k p(y=k\mid x)$
      and $\hat{\Delta}_k$ is estimated from training data.

    \column{0.44\textwidth}
      \textbf{Training-time regulariser}\\[4pt]
      Rather than correcting only at test time, ImOOD introduces a
      \emph{training loss} that directly minimises the class bias:
      \[
        \mathcal{L}_{\text{total}} =
        \mathcal{L}_{\text{CE}} + \lambda\,\mathcal{L}_{\text{bias}}
      \]
      \begin{itemize}
        \item $\mathcal{L}_{\text{bias}}$: penalises class-level gap between
              balanced and imbalanced OOD score distributions
        \item Optimised \emph{jointly} with classification
        \item \textbf{No test-time overhead}: standard scorer used as-is
        \item \textbf{Scorer-agnostic}: plug in Energy, ODIN, or Mahalanobis
      \end{itemize}
      \vspace{6pt}
      \begin{block}{Effect}
        After training, tail-class ID samples score similarly to
        head-class ID samples $\Rightarrow$ symmetric OOD detection.
      \end{block}
  \end{columns}
\end{frame}

% ----- slide: results -----
\begin{frame}{ImOOD --- Experimental Results}
  \begin{columns}[T]
    \column{0.52\textwidth}
      \textbf{AUROC on CIFAR-100-LT (IR\,=\,100), Energy scorer}
      \vspace{4pt}
      \footnotesize
      \renewcommand{\arraystretch}{1.25}
      \begin{tabular}{lcccc}
        \toprule
        Method & SVHN & LSUN & iSUN & Avg. \\
        \midrule
        Energy    & 74.3 & 76.1 & 75.8 & 75.4 \\
        ODIN      & 75.1 & 77.4 & 76.2 & 76.2 \\
        VOS       & 76.8 & 78.9 & 77.5 & 77.7 \\
        LogitAdj  & 75.9 & 77.2 & 76.0 & 76.4 \\
        GDFR      & 77.4 & 79.1 & 78.3 & 78.3 \\
        \midrule
        \textbf{ImOOD} & \good{\textbf{81.2}} & \good{\textbf{83.5}}
                       & \good{\textbf{82.7}} & \good{\textbf{82.5}} \\
        \bottomrule
      \end{tabular}
      \normalsize
      \vspace{4pt}
      {\scriptsize Refer to Table\,1 in \cite{ImOOD2024} for full results.}
      \vspace{6pt}

      \textbf{FPR95 improvement on tail classes}
      \vspace{2pt}
      \footnotesize
      \begin{tabular}{lcc}
        \toprule
        Class group & Baseline FPR95 & ImOOD FPR95 \\
        \midrule
        Head (top-20\%) & 21.4 & \neu{22.1} \\
        Mid              & 38.7 & \good{31.2} \\
        Tail (bot-20\%) & 61.3 & \good{\textbf{49.5}} \\
        \bottomrule
      \end{tabular}
      \normalsize
      \vspace{2pt}
      {\scriptsize Tail FPR95 $\downarrow$\,11.8\,pp; head FPR95 essentially unchanged.}

    \column{0.44\textwidth}
      \textbf{Key findings}
      \begin{itemize}
        \item \textbf{Tail-class FPR95 drops $\approx$12\,pp}
              on CIFAR-100-LT (IR\,=\,100)---
              largest gains exactly where class imbalance is most severe
        \item \textbf{Head-class OOD detection retained}---
              unlike rebalancing, ImOOD does not trade head
              performance for tail gains
        \item \textbf{Scorer-agnostic gains}:
              consistent improvements across Energy, ODIN, Mahalanobis,
              and KNN scorers
        \item \textbf{ID accuracy maintained}:
              $< 0.3\%$ change in classification accuracy on CIFAR
        \item \textbf{Gains scale with IR}:
              larger improvements at IR\,=\,200 than IR\,=\,10
      \end{itemize}
      \vspace{4pt}
      \begin{alertblock}{Ablation finding}
        Test-time correction alone (without $\mathcal{L}_{\text{bias}}$)
        is $\approx\!3$\,pp weaker --- training-time regularisation
        is the primary contributor.
      \end{alertblock}
  \end{columns}
\end{frame}

% =========================================================
\section{Comparative Summary \& Relevance}
% =========================================================

\begin{frame}{Comparative Summary: DCE vs.\ ImOOD}
  \footnotesize
  \renewcommand{\arraystretch}{1.4}
  \begin{tabular}{lp{5.2cm}p{5.2cm}}
    \toprule
    & \textbf{DCE} (ICML 2025) & \textbf{ImOOD} (NeurIPS 2024) \\
    \midrule
    Task
      & Domain-incremental classification
      & Out-of-distribution detection \\
    Imbalance type
      & Intra-domain class imbalance (continual)
      & Long-tailed ID training data \\
    Core mechanism
      & 3 frequency-aware experts + DES router\\
        via Gaussian pseudo-sampling
      & Class-aware bias term $\Delta_k$ +\\
        training regulariser $\mathcal{L}_{\text{bias}}$ \\
    Imbalance solution
      & Tail expert separation; pseudo-sample\\
        based routing avoids routing errors
      & Equalise OOD score distributions\\
        across head \& tail classes \\
    Overhead
      & 3$\times$ expert training + DES (lightweight)
      & One extra loss term; zero test overhead \\
    Plug-in
      & Any classifier as expert backbone
      & Any OOD scorer \\
    \midrule
    \textbf{Relevance}
      & Expert routing strategy directly applicable\\
        to MoE for network intrusion detection\\
        (addresses routing-error failure mode)
      & Per-class score calibration applicable\\
        to tail-attack detection rate improvement\\
        as a post-hoc correction \\
    \bottomrule
  \end{tabular}
  \vspace{6pt}
  \begin{block}{Shared insight}
    Both papers show that \textbf{class-specific treatment} of the imbalance
    problem---rather than global reweighting---is the key to tail-class recovery
    without degrading head-class performance.
  \end{block}
\end{frame}

% =========================================================
\section{한국어 요약 (Korean Summary)}
% =========================================================

\begin{frame}{한국어 요약}
  \begin{columns}[T]
    \column{0.48\textwidth}
      \textbf{\textcolor{dcecol}{DCE (ICML 2025)}}\\[4pt]
      DCE는 도메인 점진적 학습 환경에서 도메인 이동과 도메인 내 클래스 불균형이 복합적으로 작용하는 문제를 해결하기 위해 제안된 Mixture-of-Experts 프레임워크다.
      핵심 아이디어는 학습 클래스를 샘플 수 기준으로 헤드/균형/꼬리 세 전문가에 배정하고, 각 전문가를 독립적으로 학습시키는 것이다.
      추론 시 라우터는 실제 불균형 데이터 대신 클래스별 가우시안 분포($\mathcal{N}(\mu_k, \Sigma_k)$)에서 균등하게 생성한 의사 샘플(pseudo-samples)로 학습되므로, 꼬리 클래스에 대한 라우팅 편향이 발생하지 않는다.
      친화 행렬(affinity matrix) 기반의 소프트 결합으로 어떤 샘플도 단일 전문가에 완전히 귀속되지 않는다.
      실험에서는 4개 벤치마크 전체에서 기존 연속 학습 SOTA를 초과 달성하였으며,
      꼬리 클래스 정확도 향상이 가장 두드러졌다.
      이 프로젝트 맥락에서는 라우팅 오류 복합 문제를 직접 해소하는 설계로, MoE 구조 개선에 직접 적용 가능하다.

    \column{0.48\textwidth}
      \textbf{\textcolor{imoocol}{ImOOD (NeurIPS 2024)}}\\[4pt]
      ImOOD는 긴 꼬리 분포로 학습된 모델이 이상치(OOD) 탐지에서 보이는 클래스 비대칭 편향을 수정하는 방법이다.
      표준 OOD 탐지기는 헤드 클래스에 유리하게 편향되어, 꼬리 클래스의 정상(ID) 샘플을 OOD로 오탐하는 문제가 있다.
      ImOOD는 클래스별 편향 항 $\Delta_k = \mathbb{E}[s^{\text{bal}}|y\!=\!k] - \mathbb{E}[s^{\text{imb}}|y\!=\!k]$을 추정하고,
      이를 $\mathcal{L}_{\text{bias}}$ 형태로 학습 손실에 포함하여 훈련 단계에서 편향을 직접 줄인다.
      테스트 시에는 $S(x) = s^{\text{imb}}(x) + \hat{\Delta}_{\hat{y}}$ 로 점수를 교정하며 추가 계산 비용이 없다.
      에너지 점수·ODIN·마할라노비스 등 임의의 OOD 스코어러에 플러그인 방식으로 적용 가능하다.
      실험에서 꼬리 클래스의 FPR95가 최대 약 12\,pp 감소하고, ID 분류 정확도는 거의 유지된다.
      이 프로젝트에서는 희귀 공격 클래스 탐지율 향상을 위한 후처리 보정 기법으로 활용 가능성이 크다.
  \end{columns}
\end{frame}

% =========================================================
% References
% =========================================================
\begin{frame}[allowframebreaks]{References}
  \footnotesize
  \begin{thebibliography}{9}
    \bibitem{DCE2025}
      Yan et al.,
      \textit{Dual-Balance Collaborative Experts for Domain-Incremental Learning},
      ICML 2025. \texttt{github.com/Lain810/DCE}

    \bibitem{ImOOD2024}
      Wang et al.,
      \textit{Towards Balanced Out-of-Distribution Detection under Long-Tailed Training Data},
      NeurIPS 2024. \texttt{github.com/alibaba/imood}
  \end{thebibliography}
\end{frame}

\end{document}
